\documentclass[10pt,twocolumn]{article}
\usepackage[margin=0.75in]{geometry}
\usepackage{booktabs}
\usepackage{amsmath}
\usepackage{graphicx}
\usepackage{xcolor}
\usepackage{hyperref}
\usepackage{caption}

\definecolor{improved}{rgb}{0.0, 0.5, 0.0}
\definecolor{regressed}{rgb}{0.8, 0.0, 0.0}
\definecolor{neutral}{rgb}{0.4, 0.4, 0.4}

\title{DynLang-SLAM: Tracking Pipeline Tuning Report\\
\large Systematic Ablation Study on Replica Dataset}
\author{Ankur Guruprasad\\
EEE 515 -- Computer Vision, Arizona State University\\
\texttt{agurupr6@asu.edu}}
\date{April 2026}

\begin{document}
\maketitle

\section{Introduction}

This report documents the systematic parameter tuning and algorithmic improvements applied to the tracking pipeline of DynLang-SLAM, a 3D Gaussian Splatting-based visual SLAM system. All experiments were conducted on the Replica dataset~\cite{replica} using three scenes (\texttt{room0}, \texttt{room1}, \texttt{room2}) with 200 frames each. The primary metric is Absolute Trajectory Error (ATE RMSE) in centimeters.

\subsection{System Overview}
The tracking pipeline optimizes a 6-DOF camera pose per frame by minimizing a weighted combination of photometric (L1 + SSIM) and geometric (depth L1) losses against a differentiable Gaussian splatting map rendered via \texttt{gsplat}~\cite{gsplat}. The pose is parameterized as a quaternion (4D) + translation (3D) and optimized with Adam for a fixed number of iterations.

\subsection{Baseline Configuration}
\begin{itemize}
    \item Tracking iterations: 50
    \item Pose learning rate: 0.002 (constant)
    \item Loss weights: RGB=0.5, Depth=1.0, SSIM=0.2
    \item Render downscale: 1 (full resolution 680$\times$1200)
    \item Keyframe interval: every 5 frames
    \item Densification: gradient threshold 0.0002, interval 10
\end{itemize}

\section{Initial Baseline Results}

Table~\ref{tab:initial} shows the original ATE values before any tuning.

\begin{table}[h]
\centering
\caption{Initial baseline ATE (RMSE, cm) on Replica scenes.}
\label{tab:initial}
\begin{tabular}{lccc}
\toprule
\textbf{Scene} & \textbf{ATE (cm)} & \textbf{PSNR} & \textbf{Target} \\
\midrule
room0 & 0.72 & -- & $<$1.0 \\
room1 & 10.34 & -- & $<$5.0 \\
room2 & 18.82 & -- & $<$5.0 \\
\bottomrule
\end{tabular}
\end{table}

\section{Methods Explored}

\subsection{Depth Weight Tuning}

We varied the depth loss weight to find the optimal balance between geometric and photometric alignment.

\begin{table}[h]
\centering
\caption{Depth weight ablation (ATE RMSE in cm).}
\label{tab:depth_weight}
\begin{tabular}{lccc}
\toprule
\textbf{depth\_w} & \textbf{room0} & \textbf{room1} & \textbf{room2} \\
\midrule
1.0 (baseline) & 1.67 & 9.79 & 4.43 \\
1.2 & 1.43 & 8.15 & 3.73 \\
1.5 & -- & 9.50 & 1.77 \\
\bottomrule
\end{tabular}
\end{table}

\textbf{Finding:} \texttt{depth\_w=1.2} is the best universal compromise. Higher values (1.5) dramatically help room2 but catastrophically hurt room1, suggesting scene-dependent sensitivity to depth weighting.

\subsection{Velocity Initialization}

Constant velocity motion model: the initial pose for frame $t$ is extrapolated as $\hat{T}_t = T_{t-1} \cdot (T_{t-2}^{-1} \cdot T_{t-1})$.

\begin{table}[h]
\centering
\caption{Velocity initialization ablation.}
\label{tab:velocity}
\begin{tabular}{lcc}
\toprule
\textbf{Config} & \textbf{room0} & \textbf{room1} \\
\midrule
Without velocity init & 1.51 & 15.88 \\
With velocity init & 1.43 & 10.98 \\
\bottomrule
\end{tabular}
\end{table}

\textbf{Finding:} Velocity initialization is critical for room1 (31\% improvement) where camera motion is aggressive. Neutral for room0 where motion is smooth.

\subsection{Isotropic Regularization Weight}

Controls the strength of isotropic shape regularization on Gaussian primitives during mapping.

\begin{table}[h]
\centering
\caption{Isotropic regularization ablation on room0.}
\label{tab:iso}
\begin{tabular}{lc}
\toprule
\textbf{iso\_reg\_weight} & \textbf{room0 ATE (cm)} \\
\midrule
0.01 & 1.70 \\
0.05 & 1.43 \\
\bottomrule
\end{tabular}
\end{table}

\textbf{Finding:} Stronger isotropic regularization (\texttt{iso=0.05}) helps room0 by producing rounder, more stable Gaussians that provide a cleaner tracking signal.

\subsection{Densification}

Tested the effect of removing adaptive Gaussian densification entirely.

\begin{table}[h]
\centering
\caption{Densification ablation on room0.}
\label{tab:densify}
\begin{tabular}{lc}
\toprule
\textbf{Config} & \textbf{room0 ATE (cm)} \\
\midrule
With densification & 1.51 \\
Without densification & 2.33 \\
\bottomrule
\end{tabular}
\end{table}

\textbf{Finding:} Densification is essential -- without it, the map becomes too sparse for accurate tracking.

\subsection{Motion Magnitude Clamping}

Clamped the predicted velocity-based motion to prevent extreme extrapolations during sudden camera jerks (5$\times$ running average, 0.05m minimum).

\begin{table}[h]
\centering
\caption{Motion clamping ablation.}
\label{tab:clamp}
\begin{tabular}{lcc}
\toprule
\textbf{Config} & \textbf{room0} & \textbf{room1} \\
\midrule
Without clamping & 1.67 & 9.79 \\
With clamping & 1.78 & 7.54 \\
\bottomrule
\end{tabular}
\end{table}

\textbf{Finding:} Clamping helped room1 (23\% better) but slightly hurt room0. The tradeoff was not favorable, so this was \textbf{removed}.

\subsection{Outlier Rejection with Re-Tracking}

Detected tracking failures via loss thresholds and re-initialized from the previous frame's pose.

\begin{table}[h]
\centering
\caption{Outlier rejection ablation.}
\label{tab:outlier}
\begin{tabular}{lc}
\toprule
\textbf{Config} & \textbf{room1 ATE (cm)} \\
\midrule
Without outlier rejection & 9.79 \\
With outlier rejection & 26.14 \\
\bottomrule
\end{tabular}
\end{table}

\textbf{Finding:} \textcolor{regressed}{Catastrophic failure.} Re-tracking from the previous pose disrupted the velocity model's accumulated motion estimates, causing cascading drift. \textbf{Removed.}

\section{Failed Approaches from Literature}

Based on techniques from GS-SLAM, TVG-SLAM, CG-SLAM, and others, we implemented and tested several advanced methods. All were evaluated against the tuned baseline (\texttt{depth\_w=1.2}, velocity init, \texttt{iso=0.05}).

\subsection{Coarse-to-Fine Tracking}

Inspired by GS-SLAM~\cite{gsslam}: render at lower resolution for early iterations (coarse pose estimation), then refine at full resolution.

\begin{table}[h]
\centering
\caption{Coarse-to-fine tracking ablation on room0.}
\label{tab:c2f}
\begin{tabular}{lc}
\toprule
\textbf{Config} & \textbf{room0 ATE (cm)} \\
\midrule
Baseline (no C2F) & 1.43 \\
C2F: 4$\times$ ds, 60\% coarse & 4.07 \\
C2F: 2$\times$ ds, 40\% coarse & 3.41$^*$ \\
\bottomrule
\end{tabular}
\end{table}
{\small $^*$Combined with other changes (LR scheduling, Huber loss, depth-aware scales).}

\textbf{Finding:} \textcolor{regressed}{Significant regression.} Indoor scenes rely on high-frequency geometric details (chair legs, desk edges). Downscaling destroys the alignment signal that makes the optimizer converge. Unlike outdoor/large-scale scenes where coarse-to-fine helps, Replica rooms require full-resolution tracking throughout.

\subsection{Learning Rate Scheduling}

Applied CosineAnnealingLR to the pose optimizer, decaying from 0.002 to 0.0006 over 50 iterations.

\begin{table}[h]
\centering
\caption{LR scheduling ablation on room0.}
\label{tab:lr_sched}
\begin{tabular}{lc}
\toprule
\textbf{Config} & \textbf{room0 ATE (cm)} \\
\midrule
Constant LR (0.002) & 1.77 \\
CosineAnnealingLR + Huber & 7.03 \\
\bottomrule
\end{tabular}
\end{table}

\textbf{Finding:} \textcolor{regressed}{Severe regression.} With only 50 iterations, decaying the learning rate traps the pose in local minima before it can traverse the loss landscape. Constant LR is essential for this short optimization horizon.

\subsection{Huber Depth Loss}

Replaced L1 depth loss with Huber loss ($\delta=0.1$m) for robustness to depth outliers.

\textbf{Finding:} \textcolor{regressed}{Contributed to regression} (tested jointly with LR scheduling). Large depth errors are often the \emph{only} valid signal pulling a wrong pose back to the correct basin of attraction. Clipping these gradients (as Huber does for $|e| > \delta$) prevents convergence.

\subsection{Depth-Aware Scale Initialization}

Initialized Gaussian scales proportional to depth $\times$ pixel footprint instead of a uniform 1cm.

\textbf{Finding:} \textcolor{regressed}{Caused 30\% fewer Gaussians} (69K vs 99K). Larger initial scales meant fewer pixels were flagged as ``unmapped'' during expansion, resulting in a sparser map that hurt tracking.

\section{Successful Method: Dynamic Loss Weighting}

\subsection{Motivation}

Instead of changing the learning rate (which failed), we change the \textbf{loss weight distribution} across tracking iterations. The key insight: depth loss creates a convex basin of attraction for coarse alignment, while RGB/SSIM losses provide high-frequency signal for sub-pixel refinement.

\subsection{Method}

Given base loss weights $w_\text{rgb}$, $w_\text{depth}$, $w_\text{ssim}$, we interpolate per-iteration weights using a linear schedule parameterized by $t = i / (N-1) \in [0, 1]$:

\begin{align}
    s_\text{depth}(t) &= 2.0 - 1.5t & \text{(2.0 $\to$ 0.5)} \\
    s_\text{rgb}(t) &= 0.3 + 1.2t & \text{(0.3 $\to$ 1.5)}
\end{align}

The effective weights at iteration $i$ are:
\begin{align}
    w_\text{depth}^{(i)} &= w_\text{depth} \cdot s_\text{depth}(t) \\
    w_\text{rgb}^{(i)} &= w_\text{rgb} \cdot s_\text{rgb}(t) \\
    w_\text{ssim}^{(i)} &= w_\text{ssim} \cdot s_\text{rgb}(t)
\end{align}

This keeps the total gradient magnitude roughly constant (unlike LR scheduling) while steering the gradient \emph{direction} from geometry-dominant to appearance-dominant.

\subsection{Results}

\begin{table}[h]
\centering
\caption{Dynamic loss weighting results (ATE RMSE in cm). All scenes improved.}
\label{tab:dynamic}
\begin{tabular}{lcccc}
\toprule
\textbf{Scene} & \textbf{Original} & \textbf{Tuned BL} & \textbf{Dyn. Wt.} & \textbf{$\Delta$} \\
\midrule
room0 & 0.72 & 1.43 & \textcolor{improved}{\textbf{1.12}} & -22\% \\
room1 & 10.34 & 9.79 & \textcolor{improved}{\textbf{3.32}} & -66\% \\
room2 & 18.82 & 3.73 & \textcolor{improved}{\textbf{1.50}} & -60\% \\
\bottomrule
\end{tabular}
\end{table}

\subsection{Analysis}

Dynamic loss weighting achieves the same goal as coarse-to-fine tracking (broad convergence first, then refinement) but through the loss landscape rather than the image resolution:

\begin{itemize}
    \item \textbf{Early iterations (depth-heavy):} The depth loss provides a smooth, convex error surface that guides the optimizer into the correct basin of attraction, even when the initial pose estimate is far from ground truth.
    \item \textbf{Late iterations (RGB/SSIM-heavy):} Once the pose is approximately correct, photometric losses exploit high-frequency texture and edge information for sub-pixel alignment that depth sensors cannot provide.
\end{itemize}

The improvement is especially dramatic for room1 (66\% reduction), which features aggressive camera motion. The depth-heavy initialization phase prevents the catastrophic drift that previously occurred around frames 100--150.

\section{Other Improvements}

Two additional low-risk changes were retained:

\begin{itemize}
    \item \textbf{Quaternion normalization:} Normalizing the quaternion after each Adam step prevents the quaternion magnitude from drifting, which would otherwise distort Adam's effective learning rate.
    \item \textbf{Fast SE(3) inverse:} Replacing \texttt{torch.inverse()} with the analytical rigid-body inverse $T^{-1} = [R^\top | -R^\top t]$ in non-differentiable code paths. Mathematically identical, computationally cheaper.
\end{itemize}

\section{Summary of All Experiments}

Table~\ref{tab:summary} provides a comprehensive overview of every method tested.

\begin{table*}[t]
\centering
\caption{Complete ablation summary. All values are ATE RMSE (cm) on 200 frames. Green = improvement, Red = regression.}
\label{tab:summary}
\begin{tabular}{llcccl}
\toprule
\textbf{\#} & \textbf{Method} & \textbf{room0} & \textbf{room1} & \textbf{room2} & \textbf{Verdict} \\
\midrule
0 & Original baseline & 0.72 & 10.34 & 18.82 & Starting point \\
1 & + depth\_w=1.2 & 1.43 & 8.15 & 3.73 & \textcolor{improved}{Kept} \\
2 & + Velocity init & 1.43 & 10.98$^\dagger$ & -- & \textcolor{improved}{Kept} \\
3 & + iso\_reg=0.05 & 1.43 & -- & -- & \textcolor{improved}{Kept} \\
4 & + Motion clamping & 1.78 & 7.54 & -- & \textcolor{regressed}{Removed} \\
5 & + Outlier rejection & -- & 26.14 & -- & \textcolor{regressed}{Removed} \\
6 & + Coarse-to-fine (4$\times$) & 4.07 & -- & -- & \textcolor{regressed}{Removed} \\
7 & + All changes combined & 3.41 & -- & -- & \textcolor{regressed}{Removed} \\
8 & + LR scheduling + Huber & 7.03 & -- & -- & \textcolor{regressed}{Removed} \\
9 & + Quat norm + fast SE(3) & 1.77 & -- & -- & \textcolor{neutral}{Neutral} \\
10 & + \textbf{Dynamic loss weighting} & \textbf{1.12} & \textbf{3.32} & \textbf{1.50} & \textcolor{improved}{\textbf{Best}} \\
\bottomrule
\end{tabular}
\\[2pt]
{\small $^\dagger$Without velocity init, room1 = 15.88cm.}
\end{table*}

\section{Key Lessons Learned}

\begin{enumerate}
    \item \textbf{One change at a time.} Combining multiple modifications simultaneously caused unpredictable interactions and made debugging impossible. Sequential ablation is essential.
    \item \textbf{Pose optimization is short and non-convex.} With only 40--50 iterations, techniques that work for long training runs (LR scheduling, Huber loss) can be counterproductive.
    \item \textbf{Depth is a convex anchor.} Heavy depth weighting in early iterations creates a smooth loss landscape that prevents local minima entrapment.
    \item \textbf{Resolution matters for indoor scenes.} Coarse-to-fine tracking fails because indoor environments rely on high-frequency geometric details for alignment.
    \item \textbf{Don't fight the velocity model.} Outlier rejection with re-tracking disrupted the constant velocity assumption, causing cascading drift.
\end{enumerate}

\section{Next Steps}

\begin{itemize}
    \item Full 2000-frame evaluation on ASU Sol supercomputer (A100 GPU)
    \item Push room0 below 1.0cm target (information-driven keyframe selection)
    \item Language pipeline: CLIP + SAM + autoencoder feature fusion
    \item Dynamic object masking: YOLOv8 detection and exclusion
\end{itemize}

\begin{thebibliography}{9}
\bibitem{replica} Straub et al., ``The Replica Dataset: A Digital Replica of Indoor Spaces,'' \emph{arXiv:1906.05797}, 2019.
\bibitem{gsplat} Ye et al., ``gsplat: An Open-Source Library for Gaussian Splatting,'' \emph{arXiv:2409.06765}, 2024.
\bibitem{gsslam} Yan et al., ``GS-SLAM: Dense Visual SLAM with 3D Gaussian Splatting,'' \emph{CVPR}, 2024.
\end{thebibliography}

\end{document}
